\chapter{绪论}

\section{课题背景}
MRI 是医学教学和脑部结构观察中很常见的一类影像资料。它不使用电离辐射，对软组织的显示也比较清楚，因此常被用来观察脑部、脊柱和关节等部位。对人工智能专业的学生来说，MRI 数据也适合用来学习医学图像预处理、生成模型和三维可视化等内容。

不过，在课程实验和毕业设计中直接使用真实病例数据并不方便。医院数据通常受到隐私保护、伦理审查和数据授权限制，采集成本也比较高。公开数据集虽然能缓解这个问题，但样本来源、扫描设备和参数并不完全一致。模型训练时，这些差异会影响灰度分布和训练稳定性；课堂展示时，如果只展示一张二维切片，学习者也很难直观看到脑部结构在空间中的连续变化。

基于这些问题，本文把重点放在一个可运行的平台流程上。StyleGAN2-ADA 用于学习二维 MRI 切片的灰度和轮廓分布\cite{karras2020stylegan2,karras2020ada}，3DGS 用于把连续切片进一步组织为可展示的三维表达\cite{kerbl20233dgs}。两者在本文中的作用并不相同：StyleGAN2 负责生成二维切片，3DGS 负责三维表达和结果展示。平台则把预处理、训练、生成、点云初始化、三维训练结果和前端页面连接起来。

需要说明的是，MRI 切片不是普通相机从不同角度拍下的照片。原始 3DGS 面向多视角图像重建，而 MRI 更像是一组按层排列的灰度切片。本文没有把自然场景的 3DGS 流程直接照搬到医学图像上，而是先把连续切片转成点云，再构建规则相机参数，最后把训练日志、渲染结果和点云文件接入平台展示。这样做的目标是完成“二维生成结果如何进入三维表达”的工程验证。

\section{国内外研究现状}
医学图像生成最早常用于数据扩充。常见做法包括旋转、裁剪、缩放和灰度扰动。这些方法容易实现，但只是改变已有图像，不能真正生成新的结构。随着深度生成模型发展，GAN、变分自编码器和扩散模型开始被用于 MRI、CT、眼底图像等医学影像合成任务。

近几年的研究更重视一个问题：生成图像看起来像，并不等于一定有医学意义。Skandarani 等比较了多种 GAN 在不同医学图像上的表现，指出不同任务中的生成质量和下游效果并不一致\cite{skandarani2023gans}。Khader 等把扩散模型用于三维医学图像生成，强调合成数据需要同时考虑外观、解剖结构和切片连续性\cite{khader2023ddpm3d}。Müller-Franzes 等比较了潜在扩散模型与 GAN，也说明医学图像合成不能只看单一视觉指标\cite{mullerfranzes2023comparison}。Akbar 等在脑肿瘤 MR 合成图像上进行分割实验，进一步说明合成数据需要放到具体任务中验证\cite{akbar2024syntheticmr}。这些研究给本文一个直接启发：本平台的生成结果适合教学展示和工程实验，不应被解释为临床诊断数据。

在二维生成方面，GAN 为本文提供了基本方法框架\cite{goodfellow2014gan}。StyleGAN 通过风格控制提高图像合成质量\cite{karras2019stylegan}，StyleGAN2 进一步减少伪影并改进训练稳定性\cite{karras2020stylegan2}。StyleGAN2-ADA 针对有限数据训练加入自适应增强机制\cite{karras2020ada}。本文使用的 IXI 脑部 MRI 切片数量虽然不少，但病例来源仍有限，且图像主要是灰度结构。因此，选择 StyleGAN2-ADA 更符合当前项目的数据条件和工程实现难度。

在三维表达方面，传统医学可视化常使用体绘制、表面重建或体素显示，这些方法更适合直接展示真实体数据。本文还需要处理生成切片，因此要先解决二维结果如何进入三维空间的问题。3DGS 使用点云和三维高斯表达空间结构\cite{kerbl20233dgs}。已有研究把 3DGS 用于内窥镜、手术场景和医学相关三维视觉任务\cite{bonilla2024gaussianpancakes,huang2024endo4dgs,guo2024freesurgs,zeng2024surgical3dgs}。这些工作说明 3DGS 具有可扩展性，但 MRI 切片序列与多视角照片仍有明显差别。本文的工作重点正是在这个差别上做工程适配：把连续切片、初始点云和规则相机参数组织成可训练、可展示的输入。

\section{项目研究意义}
本课题的意义主要体现在工程整合和教学展示上。过去做类似实验时，数据预处理、StyleGAN2 训练、切片生成、点云构建和三维展示往往分散在多个脚本和文件夹中。本文把这些步骤放进同一条流程，用户可以通过页面查看数据、生成结果、任务日志和三维资产，实验过程更容易复核。

从医学图像智能处理角度看，本文尝试把二维 MRI 切片生成和三维表达连接起来。StyleGAN2-ADA 负责生成切片，连续切片序列负责建立层间关系，点云和 3DGS 则负责把这些二维结果组织成可展示的空间形式。这样处理后，生成结果不只是一张图片，而是可以继续进入三维展示流程。

从教学应用角度看，FastAPI + React 平台降低了操作门槛。用户不用记住复杂命令，也不用在训练目录中手动查找结果文件。平台可以提供参数输入、运行状态、图像预览、点云展示和风险提示。对课程演示、实验复核和毕业答辩来说，这种方式比单独展示脚本输出更直观。

\section{论文主要内容与章节安排}
围绕上述目标，论文工作主要包括四个部分。

第一部分是 MRI 数据预处理。系统读取 IXI 数据集中 NIfTI 格式的 T1 脑部 MRI 数据，提取轴向切片，并完成百分位归一化、尺寸统一和低信息切片过滤。处理后的数据可用于 StyleGAN2-ADA 训练。

第二部分是二维 MRI 切片生成。系统基于预处理切片训练 StyleGAN2-ADA，并在平台中支持权重加载、随机种子设置、截断参数调节和批量切片生成。

第三部分是面向三维表达的序列组织和 3DGS 训练。系统通过潜空间插值得到连续切片，再依据像素坐标和灰度值写出 PLY 点云。结合规则相机参数后，系统完成 3DGS 训练，并保存 checkpoint、训练后点云和渲染图像。

第四部分是 FastAPI + React 交互平台。后端提供生成、分析、任务、资产、体数据和 AI 助手等接口；前端提供仪表盘、生成工作台、医学影像查看器、报告、可信性评估、点云展示和教学中心等页面。用户可以在统一界面中完成主要实验操作。

本文结构安排如下。第 1 章介绍研究背景、研究现状、研究意义和论文结构。第 2 章介绍 MRI、NIfTI、GAN、StyleGAN2、3DGS、FastAPI、React 和 Vite 等相关技术。第 3 章进行系统需求分析。第 4 章给出系统总体设计。第 5 章说明系统关键模块实现。第 6 章进行测试与结果分析。第 7 章说明本文工作完成情况与后续改进方向。
